% =============== Discussion ===================================================
% Goal: 1-1.5 pages
% =============================================================================

\subsection{Anchor Design as the Dominant Variable}
\label{sec:discussion-anchors}

The single most consequential lesson from this work is that
\emph{anchor design dominates technical sophistication}. The
+0.42 AUC swing from changing how anchors are written (Step 1)
dwarfs the gains from switching embedders (Step 3, +0.04) or
adding hard examples (Step 4, +0.09). For practitioners building
similar anchor-based detection systems, this suggests budget
priorities skewed toward anchor curation rather than model
selection.

The principle generalizes. Drift detection at the prompt-text
layer is fundamentally a similarity-search problem; cosine
similarity is sensitive to grammatical mood, lexical distribution,
and topic distribution. Anchors that systematically differ from
real prompts in any of these axes will yield poor detection
regardless of how strong the underlying embedder is.

\subsection{Black-box vs.\ White-box Persona Monitoring}
\label{sec:discussion-bw}

\citet{personavectors2025}'s white-box approach has one major
advantage and one major disadvantage relative to ours.

\paragraph{Advantage of white-box.} Activation projections operate
on the LLM's actual computational state and can detect drift even
when the surface-level prompt looks innocuous. A prompt designed
to elicit sycophancy without explicit sycophantic surface forms
may activate sycophancy directions in the residual stream while
appearing aligned to a black-box anchor matcher. This is the limit
case of black-box detection.

\paragraph{Disadvantage of white-box.} The method requires model
weights and (typically) a training set of contrastive prompts to
extract the persona vectors. End-users of proprietary LLMs
(Claude, GPT-4, Gemini) have neither.

\paragraph{Complementarity.} A robust persona monitoring stack
plausibly combines both: white-box detection where model access
is available (model trainers, alignment researchers, organizations
with self-hosted Llama/Mistral deployments), and black-box
detection in user-space hooks for the long tail of users
interacting with closed APIs. Our work targets the latter.

\subsection{Why Reranker Helps Retrieval but Not Drift}
\label{sec:discussion-rerank}

The asymmetric finding---bge-reranker-v2-m3 gives +0.105 MRR on
LongMemEval-S retrieval but +0.0008 AUC on drift
(Section~\ref{sec:eval-rerank-null})---is informative.

In retrieval, candidates are full conversational sessions with
abundant token-level interaction signal: answer terms appearing
near question terms, paraphrastic framing, dependency relations.
A cross-encoder can exploit this to break ties and re-rank
candidates the bi-encoder gets wrong.

In drift detection, anchors are short imperative sentences. The
``token interaction'' between a prompt and an anchor is largely
keyword overlap, which a bi-encoder already captures.
Cross-encoders cannot manufacture signal that does not exist in
the data.

\subsection{The Subset 4 vs.\ Subset 12 Lesson}
\label{sec:discussion-subset}

A non-technical lesson worth recording: during development we ran
the reranker eval on subset 4 (one question per type for the four
mainstream types). Subset 4 yielded $\Delta$ MRR $+0.001$, leading
us to a near-published null conclusion. Increasing to subset 12
(two per type, including the hardest types) revealed
$\Delta$ MRR $+0.105$. The lift signal was concentrated in
question types absent from subset 4.

The takeaway: \emph{do not draw conclusions from $n=4$ retrieval
benchmarks}, especially when the dataset has known per-type variance
\citep{longmemeval2024}. We were initially confident in the null
because the experimental setup was correct (clean control,
deterministic pipeline); the failure was statistical rather than
methodological.

\subsection{Active Learning in Practice}
\label{sec:discussion-al}

Our active learning module surfaces alerts whose drift score is
near the decision boundary ($|\mathrm{drift}| < 0.05$) for
preferential user labeling. The theoretical case is straightforward:
labels near the boundary disambiguate the largest amount of
uncertainty. In practice, the long-term value of this loop depends
on user willingness to label, which we have not yet measured at
scale (current evaluation is on synthetic alerts seeded into the
log file).

We expect the labeling burden to follow a power law: the first 10
labels yield large gains; the next 100 yield diminishing
incremental AUC; beyond 1000 labels the marginal gain approaches
zero. Verifying this empirically is part of the long-term
deployment study.

\subsection{Implications for Production Deployment}
\label{sec:discussion-prod}

Three practical recommendations emerge from our experience:

\textbf{(1) Tune anchors per-domain.} The 25+35-anchor default set
shipped with nautilus-compass v0.7.1 targets a Chinese/English mixed
engineering context. Deployment in legal, medical, or financial
domains will likely show degraded AUC unless anchors are
re-authored or replaced from one of our domain starter packs
(legal/medical/finance).

\textbf{(2) Filter system-injected prompts.} In a real coding
agent, hook input is not exclusively user-issued; harness-injected
notifications, monitor outputs, and tool reminders share the same
input channel. We observed $\sim 80\%$ of false alerts in our own
session traces traced to such injection. A simple regex filter
(Section~\ref{sec:method-fpfilter}) addresses this.

\textbf{(3) Combine drift detection with retrieval.} Drift
detection alone tells the agent ``you may be about to do
something bad''; retrieval tells it ``here's the relevant prior
guidance''. Both layered into the same hook output gives the LLM
the best chance of self-correcting.
